\section{The IGE Rises With Age, and What Rises}\label{sec:ageprofile}

The IGE of sons' earnings rises by about three fifths between the late twenties and the forties. In the PSID's
baseline sample it is 0.276$^{***}$ at 25--29 and 0.472$^{***}$ at 40--44, and it stays near that level to 55
(0.455$^{***}$ at 50--55) (Table~\ref{tab:ige_baseline}). The
same sons followed from 25 to 49 show the same rise, and so do two other cohorts in the NLSY79 and the NLSY97. Ranks of
family income rise only into the early thirties, so the 31--37 window of \citet{chetty2020race} lies on the plateau of the
rank profile and before the end of the log profile's rise. What rises is mostly the pay-off of what sons bring to the
labor market---their skill at entry---while the parents' wage and wealth carry most of the level; one thing parents do
carries a share that rises with the son's age, the father's layoffs, which the rest of the paper is about. The rise is
White: the IGE of Black sons is flat in age.

\subsection{The profile}

Pooled across races and weighted, the IGE of sons' own earnings is 0.276$^{***}$ (0.022) at 25--29, 0.403$^{***}$ at
30--34, 0.430$^{***}$ at 35--39, 0.472$^{***}$ at 40--44 and 0.443$^{***}$ at 45--49, a change of 0.167$^{***}$ ($p<0.001$)
between the first and the fifth band; that of their family income goes from 0.326$^{***}$ (0.023) to 0.521$^{***}$ at
40--44. Equality across bands is rejected ($\chi^2(5) = 41.3$, $p<0.001$) and the trend is 0.059$^{***}$ (0.022) per ten
years of age.\footnote{The baseline sample of the paper: 4,770 sons with parental income over at least three of their
ages 13--19 (Section~\ref{sec:data}); band regressions with single-age dummies, PSID weights, standard errors clustered
by 1968 family. On all sons with a measured parental income the profile is the same at 25--29 (0.276) and somewhat
higher later (0.492 at 40--44; Appendix Table~\ref{tab:ige_sons}). Code: \texttt{113.BaselineProfile.py},
\texttt{26.IGEByAge.py}.}

\begin{table}[!htbp]
\centering
\caption{The IGE of sons' earnings by age at measurement, baseline sample (PSID)}
\label{tab:ige_baseline}
\resizebox{\ifdim\width>\textwidth\textwidth\else\width\fi}{!}{%
\resizebox*{!}{\ifdim\totalheight>0.85\textheight 0.85\textheight\else\totalheight\fi}{%
\begin{tabular}{lccccc}
\toprule
 & \begin{tabular}[b]{@{}c@{}}Pooled\end{tabular} & \begin{tabular}[b]{@{}c@{}}White\end{tabular} & \begin{tabular}[b]{@{}c@{}}Black\end{tabular} & \begin{tabular}[b]{@{}c@{}}White $-$ Black\end{tabular} & \begin{tabular}[b]{@{}c@{}}Pooled, family income\end{tabular} \\
 & (1) & (2) & (3) & (4) & (5) \\
\midrule
\multicolumn{6}{l}{\textit{A. IGE of ln earnings by age band (column 5: ln family income)}} \\
25--29 & $0.276^{***}$ & $0.218^{***}$ & $0.242^{***}$ & $-0.024$ & $0.326^{***}$ \\
 & (0.022) & (0.025) & (0.031) & (0.040) & (0.023) \\
30--34 & $0.403^{***}$ & $0.371^{***}$ & $0.264^{***}$ & $0.106^{**}$ & $0.460^{***}$ \\
 & (0.026) & (0.029) & (0.031) & (0.042) & (0.027) \\
35--39 & $0.430^{***}$ & $0.419^{***}$ & $0.257^{***}$ & $0.161^{***}$ & $0.482^{***}$ \\
 & (0.030) & (0.034) & (0.038) & (0.051) & (0.032) \\
40--44 & $0.472^{***}$ & $0.485^{***}$ & $0.179^{***}$ & $0.306^{***}$ & $0.521^{***}$ \\
 & (0.035) & (0.041) & (0.045) & (0.061) & (0.035) \\
45--49 & $0.443^{***}$ & $0.457^{***}$ & $0.185^{***}$ & $0.272^{***}$ & $0.489^{***}$ \\
 & (0.048) & (0.055) & (0.057) & (0.079) & (0.048) \\
50--55 & $0.455^{***}$ & $0.441^{***}$ & $0.218^{***}$ & $0.223^{**}$ & $0.546^{***}$ \\
 & (0.063) & (0.069) & (0.070) & (0.098) & (0.057) \\
Change 45--49 minus 25--29 & $0.167^{***}$ &  &  &  & $0.164^{***}$ \\
 & (0.048) &  &  &  & (0.045) \\
Sons at 25--29 / 45--49 & 4,053 / 1,483 &  &  &  &  \\
\addlinespace
\multicolumn{6}{l}{\textit{B. Change 45--49 minus 25--29 in the pooled earnings IGE, by cohort and net of calendar time}} \\
Born 1951--64 & $0.162^{***}$ &  &  &  &  \\
 & (0.054) &  &  &  &  \\
Born 1965--79 & $0.160^{**}$ &  &  &  &  \\
 & (0.074) &  &  &  &  \\
Born 1980--89 (to 40--44) & $0.142^{*}$ &  &  &  &  \\
 & (0.080) &  &  &  &  \\
Income-year fixed effects & $0.149^{***}$ &  &  &  &  \\
 & (0.047) &  &  &  &  \\
Income years 1990--2010 only & $0.198^{***}$ &  &  &  &  \\
 & (0.058) &  &  &  &  \\
\bottomrule
\end{tabular}}}
\par\smallskip
\begin{minipage}{\textwidth}\footnotesize \textit{Notes:} The paper's baseline sample: 4,770 sons with parental income over at least three of their ages 13--19, son-years at 25--55 (earnings: valid years). Band regressions with single-age dummies; pooled estimates use PSID weights, race-specific ones are unweighted; standard errors clustered by 1968 family. Appendix Table~\ref{tab:ige_sons} reports the profile on all sons with a measured parental income. Source: \texttt{113.BaselineProfile.py}; table built by \texttt{100.PaperBFableTables.py}. $^{***}$ $p<0.01$, $^{**}$ $p<0.05$, $^{*}$ $p<0.1$.\end{minipage}
\end{table}


The rise is not a change in who is observed. In the panel of sons observed in every band of 25--49 the earnings IGE is
0.272$^{***}$ at 25--29, 0.386$^{***}$ at 30--34, 0.432$^{***}$ at 35--39, 0.473$^{***}$ at 40--44 and 0.478$^{***}$ at
45--49, a trend of 0.100$^{***}$ (0.027) per ten years. It is the intergenerational face of the life-cycle bias of
\citet{haider2006life}: following the same sons from 25--29 to 35--39, their family income grows faster with parental
income by 0.210$^{***}$ (0.031), whereas between 30--34 and 40--44 the growth gradient is 0.021 (0.036), so the rise is
concentrated before the mid-thirties.\footnote{Appendix Table~\ref{tab:growth}; \texttt{output/ige\_by\_age.md},
Sections 1b and 5.1.} Nor is it calendar time. The rise from 25--29 to 45--49 is 0.162$^{***}$ among sons born 1951--64,
0.160$^{**}$ among those born 1965--79 and 0.142$^{*}$ to 40--44 among those born 1980--89, cohorts observed at the same
ages in different decades; it is 0.149$^{***}$ with income-year fixed effects in every band and 0.198$^{***}$ in a fixed
window of income years 1990--2010, in which the bands hold different cohorts (panel B of Table~\ref{tab:ige_baseline}).
The rise replicates in the NLSY79, where the earnings IGE of sons born 1957--64 rises from 0.340$^{***}$ at 25--29 to
0.506$^{***}$ at 50--54 (an opening of 0.166$^{***}$, $p<0.001$), and in the NLSY97 (a trend of 0.063$^{***}$ over
25--44). Daughters' family income shows the same rise (Appendix Tables~\ref{tab:ige_daughters} and \ref{tab:nlsyage}).

\begin{table}[!htbp]
\centering
\caption{Age or period? The pooled rise of the IGE and the son's AFQT in two cohorts}
\label{tab:age_period}
\resizebox{\ifdim\width>\textwidth\textwidth\else\width\fi}{!}{%
\resizebox*{!}{\ifdim\totalheight>0.85\textheight 0.85\textheight\else\totalheight\fi}{%
\begin{tabular}{lccc}
\toprule
 & \begin{tabular}[b]{@{}c@{}}25--29\end{tabular} & \begin{tabular}[b]{@{}c@{}}45--49\end{tabular} & \begin{tabular}[b]{@{}c@{}}Rise\end{tabular} \\
 & (1) & (2) & (3) \\
\midrule
\multicolumn{4}{l}{\textit{A. PSID sons, ln earnings, IGE by band}} \\
Baseline & $0.283^{***}$ & $0.492^{***}$ & $0.209^{***}$ \\
 & (0.025) & (0.052) & (0.050) \\
Born 1951--64 & $0.296^{***}$ & $0.512^{***}$ & $0.216^{***}$ \\
 & (0.035) & (0.061) & (0.060) \\
Born 1965--79 & $0.299^{***}$ & $0.446^{***}$ & $0.147^{*}$ \\
 & (0.051) & (0.082) & (0.076) \\
Born 1980--89 (to 40--44) & $0.290^{***}$ & $0.455^{***}$ & $0.165^{*}$ \\
 & (0.045) & (0.085) & (0.086) \\
Income-year fixed effects & $0.298^{***}$ & $0.490^{***}$ & $0.192^{***}$ \\
 & (0.024) & (0.053) & (0.050) \\
IGE $\times$ cohort, at the 1965 cohort & $0.303^{***}$ & $0.491^{***}$ & $0.187^{***}$ \\
 & (0.027) & (0.053) & (0.050) \\
Period-varying IGE, at 1995 & $0.294^{***}$ & $0.593^{***}$ & $0.300^{***}$ \\
 & (0.025) & (0.084) & (0.084) \\
Income years 1990--2010 only & $0.286^{***}$ & $0.525^{***}$ & $0.240^{***}$ \\
 & (0.038) & (0.058) & (0.066) \\
\addlinespace
\multicolumn{4}{l}{\textit{B. The son's AFQT in the IGE of ln earnings (columns 1--2: 25--29 and 40--44)}} \\
NLSY97 (born 1980--84): IGE & $0.175^{***}$ & $0.258^{***}$ &  \\
 & (0.021) & (0.027) &  \\
NLSY97: son's AFQT contribution & $0.033^{***}$ & $0.054^{***}$ &  \\
 & (0.006) & (0.010) &  \\
NLSY79 (born 1957--64): son's AFQT contribution & $0.077^{***}$ & $0.131^{***}$ &  \\
 & (0.010) & (0.014) &  \\
\bottomrule
\end{tabular}}}
\par\smallskip
\begin{minipage}{\textwidth}\footnotesize \textit{Notes:} Panel A: Black and White sons pooled with PSID weights, ages 25--49; single-age dummies and the parent's age quartic; rise = 45--49 minus 25--29 (40--44 for the 1980--89 cohort); standard errors clustered by 1968 family. Panel B: exact Gelbach split with the son's AFQT as its own group (NLSY97: ASVAB math-verbal percentile, standardised within birth year; parental income = the 1996 parent report; NLSY79: AFQT-3, 1979--81 parental income); 200 cluster-bootstrap draws of households. Source: \texttt{95.PooledRiseChecks.py}, \texttt{94.NLSY97AFQT.py}, \texttt{87.NLSY79YouthAFQT.py}; table built by \texttt{100.PaperBFableTables.py}. $^{***}$ $p<0.01$, $^{**}$ $p<0.05$, $^{*}$ $p<0.1$.\end{minipage}
\end{table}


\subsection{Ranks}

Ranks rise less and earlier. The rank--rank slope of sons' family income is 0.328$^{***}$ (0.020) at 25--29 and
0.427$^{***}$ at 30--34, and then flat: 0.436$^{***}$ at 35--39, 0.423$^{***}$ at 40--44 and 0.380$^{***}$ at 45--49
(Appendix Table~\ref{tab:rankage}). The PSID slope in the 31--37 window, 0.442$^{***}$, lies on the plateau. Logs and
ranks differ because the log IGE also registers the dispersion of sons' log income relative to their parents', which
grows with age and which ranks normalise away.\footnote{Within the same children and weights, the widening of the
log profile is mostly the ratio of the standard deviations of children's and parents' log income;
\texttt{output/ige\_by\_age.md}, Section 5.4.} An IGE measured at 31--37 therefore captures the whole rise in ranks and
stops before the log IGE reaches its level of the forties.

